\documentclass[12pt,a4paper]{article}

% Paket
\usepackage[utf8]{inputenc}
\usepackage[swedish]{babel}
\usepackage[margin=2.5cm]{geometry}
\usepackage{graphicx}
\usepackage{xcolor}
\usepackage{hyperref}
\usepackage{listings}
\usepackage{tikz}
\usepackage{enumitem}
\usepackage{booktabs}
\usepackage{array}
\usepackage{verbatim}

% Färger
\definecolor{brandgreen}{RGB}{0,112,74}
\definecolor{lightgray}{RGB}{245,245,245}
\definecolor{codegray}{RGB}{50,50,50}

% Hyperref inställningar
\hypersetup{
  colorlinks=true,
  linkcolor=brandgreen,
  urlcolor=brandgreen,
  citecolor=brandgreen
}

% Listings för kodexempel
\lstset{
  basicstyle=\small\ttfamily,
  backgroundcolor=\color{lightgray},
  breaklines=true,
  frame=single,
  numbers=left,
  numberstyle=\tiny\color{codegray}
}

% Titel
\title{\textbf{API Endpoints Specifikation} \\ \large Komplett Backend-mappning för Onboarding-flödet}
\author{Celestial AB - Roaring}
\date{\today}

\begin{document}

\maketitle
\tableofcontents
\newpage

\section{Introduktion}

Detta dokument beskriver \textbf{alla API endpoints} som behövs för att implementera den kompletta onboarding-processen enligt \texttt{Onboardin\_app\_ny.tex}.

\subsection{Dokumentstruktur}

Dokumentet är uppdelat i två huvuddelar:

\textbf{Del 1: Autentiseringssidor (Utan Sidebar)} -- Sidor som visas innan användaren är inloggad, inklusive landing page, registrering, verifiering, inloggning och lösenordsåterställning. Dessa sidor har ingen sidebar.

\textbf{Del 2: Onboarding Slides (Med Sidebar)} -- Slides som visas efter inloggning och utgör själva onboarding-processen. Dessa har sidebar med progressindikator och navigering.

För varje slide/steg i onboarding-flödet dokumenteras:
\begin{itemize}[noitemsep]
  \item Vilka API endpoints som behövs
  \item HTTP-metod (GET, POST, PUT, DELETE)
  \item Request-format (JSON schema, multipart/form-data för filer)
  \item Response-format (success och error cases)
  \item Affärslogik och valideringar
  \item Externa integrationer (Bolagsverket, Skatteverket, BankID, etc.)
  \item Databastabeller som påverkas
\end{itemize}

\subsection{Teknisk Stack}
\begin{itemize}[noitemsep]
  \item \textbf{Backend:} FastAPI (Python 3.11+)
  \item \textbf{Databas:} PostgreSQL 15+
  \item \textbf{Authentication:} JWT tokens (access + refresh)
  \item \textbf{File Storage:} AWS S3 eller lokal disk under utveckling
  \item \textbf{OCR:} Google Cloud Vision API eller Tesseract
  \item \textbf{Externa API:er:} Bolagsverket, Skatteverket, BankID, SendGrid
\end{itemize}

\subsection{Base URL}
\begin{verbatim}
Development: http://localhost:8000/api
Production:  https://api.roaring.se/api
\end{verbatim}

\subsection{Autentisering}
Alla endpoints (utom registrering, login, verify) kräver JWT token i header:
\begin{verbatim}
Authorization: Bearer <access_token>
\end{verbatim}

\newpage

% ############################################################################
% ############################################################################
% DEL 1: AUTENTISERINGSSIDOR (UTAN SIDEBAR)
% ############################################################################
% ############################################################################

% ============================================================================
% SIDA 1: LANDING PAGE (Hero-sektion)
% ============================================================================

\section{Landing Page (Hero-sektion)}

\subsection{API-endpoints Sammanfattning}
\begin{itemize}[noitemsep]
  \item \texttt{GET /api/auth/session} — Kontrollera om användaren är inloggad (körs automatiskt om JWT finns i localStorage)
  \item \texttt{POST /api/auth/refresh} — Förnya access token om den har upphört
  \item \textbf{Antal anrop vid sidladdning:} 0–2 (beroende på om tokens finns i localStorage)
\end{itemize}

\subsection{Översikt}
Landing page är den första sidan användaren ser. Den innehåller statiskt innehåll (text, features-lista, CTA-knappar) och kräver \textbf{inga API-anrop vid laddning}.

\subsection{Funktionalitet}
\begin{itemize}[noitemsep]
  \item Visa företagsinformation och värdeproposition
  \item Knappar: "Logga in", "Registrera", "Börja onboarding nu", "Se demo"
  \item Alla knappar gör frontend-routing (ingen backend-kommunikation)
  \item \textbf{Kontrollera localStorage} för sparade JWT tokens vid sidladdning
\end{itemize}

\subsection{localStorage Check och Auth-flöde}
När React-appen laddas (i \texttt{App.jsx} eller auth context) kontrolleras localStorage för JWT tokens:

\textbf{Steg 1: Kolla localStorage}
\begin{verbatim}
const accessToken = localStorage.getItem('accessToken');
const refreshToken = localStorage.getItem('refreshToken');
\end{verbatim}

\textbf{Steg 2: Om tokens finns}
\begin{itemize}[noitemsep]
  \item Anropa \texttt{GET /api/auth/session} med \texttt{Authorization: Bearer <accessToken>}
  \item Om \texttt{200 OK} och \texttt{authenticated: true} → Användaren är inloggad
  \item Om \texttt{401 Unauthorized} → Access token har upphört, försök refresh
\end{itemize}

\textbf{Steg 3: Om access token upphört}
\begin{itemize}[noitemsep]
  \item Anropa \texttt{POST /api/auth/refresh} med refresh token
  \item Om \texttt{200 OK} → Spara ny access token, fortsätt som inloggad
  \item Om \texttt{401} → Refresh token upphört, logga ut (rensa localStorage)
\end{itemize}

\textbf{Steg 4: Om inga tokens finns}
\begin{itemize}[noitemsep]
  \item Användaren är utloggad, visa "Logga in" + "Registrera" knappar
\end{itemize}

\textbf{Knappen "Börja onboarding nu":}
\begin{itemize}[noitemsep]
  \item Om användaren HAR giltiga tokens → Frontend-routing till \texttt{/inledning} (ingen API-anrop)
  \item Om användaren INTE har tokens → Frontend-routing till \texttt{/login?redirect=/inledning}
  \item Inget dedikerat API-endpoint behövs, bara conditional routing i React
\end{itemize}

\subsection{Endpoint 1: GET /api/auth/session}

\subsubsection{Beskrivning}
Kontrollera om användaren har en aktiv session för att visa rätt UI (inloggad vs utloggad). Anropas automatiskt om JWT token finns i localStorage.

\subsubsection{Request}
\begin{itemize}[noitemsep]
  \item \textbf{Headers:} \texttt{Authorization: Bearer <token>} (optional)
  \item \textbf{Body:} None
\end{itemize}

\subsubsection{Response - Success (200 OK - Authenticated)}
\begin{lstlisting}
{
  "authenticated": true,
  "user": {
    "id": "550e8400-e29b-41d4-a716-446655440000",
    "email": "chef@revisionstockholm.se",
    "role": "byrachef",
    "firm": {
      "id": "660e8400-e29b-41d4-a716-446655440001",
      "name": "Revision Stockholm AB",
      "orgNr": "556123-4567"
    }
  }
}
\end{lstlisting}

\subsubsection{Response - Not Authenticated (200 OK)}
\begin{lstlisting}
{
  "authenticated": false
}
\end{lstlisting}

\subsubsection{Implementation Notes}
\begin{itemize}[noitemsep]
  \item Körs automatiskt om access token finns i localStorage
  \item Token valideras mot JWT secret och expiry time
  \item Om valid → Query databas för user + firm info
  \item Om invalid → Returnera \texttt{authenticated: false} (ingen DB query)
  \item Frontend uppdaterar UI baserat på response
\end{itemize}

\subsection{Endpoint 2: POST /api/auth/refresh}

\subsubsection{Beskrivning}
Förnya access token med refresh token om access token har upphört.

\subsubsection{Request}
\begin{lstlisting}
{
  "refreshToken": "eyJhbGciOiJIUzI1NiIsInR5cCI6IkpXVCJ9..."
}
\end{lstlisting}

\subsubsection{Response - Success (200 OK)}
\begin{lstlisting}
{
  "accessToken": "eyJhbGciOiJIUzI1NiIsInR5cCI6IkpXVCJ9...",
  "expiresIn": 3600
}
\end{lstlisting}

\subsubsection{Response - Error (401 Unauthorized)}
\begin{lstlisting}
{
  "error": "invalid_token",
  "message": "Refresh token har upphort eller ar ogiltig"
}
\end{lstlisting}

\subsubsection{Implementation Notes}
\begin{itemize}[noitemsep]
  \item Refresh token valideras mot databas (tabell: \texttt{refresh\_tokens})
  \item Kontrollera expiry time (vanligtvis 30 dagar)
  \item Om valid → Generera ny access token (15 min expiry)
  \item Om invalid → Användaren måste logga in igen
  \item Frontend sparar ny access token i localStorage
\end{itemize}

\subsection{Sammanfattning}
\begin{itemize}[noitemsep]
  \item \textbf{localStorage undersöks} vid sidladdning efter JWT tokens
  \item \textbf{Om tokens finns:} 1-2 API-anrop (session check + ev. refresh)
  \item \textbf{Om inga tokens:} 0 API-anrop
  \item \textbf{Endpoints:} \texttt{GET /api/auth/session} + \texttt{POST /api/auth/refresh}
  \item Alla knappar gör frontend-routing (inga nya API-anrop)
\end{itemize}

\newpage




% ============================================================================
% SIDA 2: REGISTRERING (Autentiseringssida utan sidebar)
% ============================================================================

\section{Registrering (Sida 2)}

\subsection{API-endpoints Sammanfattning}
\begin{itemize}[noitemsep]
  \item \texttt{POST /api/auth/register} — Skapa nytt konto med email/lösenord, skicka 6-siffrig kod
  \item \texttt{POST /api/auth/google} — Registrera eller logga in via Google OAuth
  \item \texttt{GET /api/auth/check-email} — Real-time validering om email redan är registrerad
  \item \textbf{Externa integrationer:} Cloudflare Turnstile (bot-skydd), SendGrid (email), Google OAuth
\end{itemize}

\subsection{Översikt}
Registreringssidan låter användaren skapa ett nytt konto med email + lösenord eller via Google OAuth. Ett bot-skydd (Cloudflare Turnstile) används för att förhindra automatiserad spam.

\subsection{Funktionalitet}
\begin{itemize}[noitemsep]
  \item Formulär: Email, Lösenord, Bekräfta lösenord
  \item Lösenordskrav: Minst 12 tecken, en siffra, en versal
  \item Cloudflare Turnstile checkbox (bot-skydd)
  \item Alternativ: "Logga in med Google" (OAuth)
  \item Real-time validering av email (om redan registrerad)
  \item Vid lyckad registrering → Redirect till \texttt{/verify} med email i state
\end{itemize}

\subsection{Endpoint 1: POST /api/auth/register}

\subsubsection{Beskrivning}
Skapa nytt användarkonto, generera 6-siffrig verifieringskod, skicka via email.

\subsubsection{Request}
\begin{lstlisting}
{
  "email": "chef@revisionstockholm.se",
  "password": "SecurePass123!",
  "passwordConfirm": "SecurePass123!",
  "captchaToken": "cloudflare-turnstile-response-token"
}
\end{lstlisting}

\subsubsection{Response - Success (201 Created)}
\begin{lstlisting}
{
  "success": true,
  "userId": "550e8400-e29b-41d4-a716-446655440000",
  "message": "Konto skapat. En 6-siffrig verifieringskod har skickats till din e-post.",
  "email": "chef@revisionstockholm.se"
}
\end{lstlisting}

\subsubsection{Response - Error (400 Bad Request)}
\textbf{Email redan registrerad:}
\begin{lstlisting}
{
  "error": "email_exists",
  "message": "Email-adressen ar redan registrerad"
}
\end{lstlisting}

\textbf{Svagt lösenord:}
\begin{lstlisting}
{
  "error": "weak_password",
  "message": "Losenordet maste vara minst 12 tecken, innehalla en siffra och en versal"
}
\end{lstlisting}

\textbf{Captcha misslyckades:}
\begin{lstlisting}
{
  "error": "captcha_failed",
  "message": "Bot-skyddet misslyckades. Forsom igen."
}
\end{lstlisting}

\subsubsection{Implementation Notes}
\begin{itemize}[noitemsep]
  \item Validera email format (regex)
  \item Kontrollera att email inte finns i \texttt{users} tabell
  \item Validera lösenord: minst 12 tecken, \texttt{/[A-Z]/}, \texttt{/[0-9]/}
  \item Kontrollera \texttt{password === passwordConfirm}
  \item Validera Cloudflare Turnstile token mot Cloudflare API
  \item Hasha lösenord med bcrypt (cost factor 12)
  \item Skapa user i databas med \texttt{verified: false}
  \item Generera 6-siffrig kod (t.ex. \texttt{482916})
  \item Spara kod i \texttt{verification\_codes} tabell med 15 min expiry
  \item Skicka email via SendGrid: \textit{"Din verifieringskod är: 482916. Koden är giltig i 15 minuter."}
  \item Redirect frontend till \texttt{/verify} med email i state
\end{itemize}

\subsubsection{Databastabeller}
\begin{itemize}[noitemsep]
  \item \texttt{users} (id, email, password\_hash, verified, created\_at, updated\_at)
  \item \texttt{verification\_codes} (id, email, code, expires\_at, attempts, created\_at)
\end{itemize}

\subsection{Endpoint 2: POST /api/auth/google}

\subsubsection{Beskrivning}
Logga in eller registrera med Google OAuth. Ingen email-verifiering behövs (Google har redan verifierat).

\subsubsection{Request}
\begin{lstlisting}
{
  "code": "google-oauth-authorization-code",
  "redirectUri": "http://localhost:5173/auth/callback"
}
\end{lstlisting}

\subsubsection{Response - Success (200 OK för befintlig, 201 Created för ny användare)}
\begin{lstlisting}
{
  "success": true,
  "accessToken": "eyJhbGciOiJIUzI1NiIsInR5cCI6IkpXVCJ9...",
  "refreshToken": "eyJhbGciOiJIUzI1NiIsInR5cCI6IkpXVCJ9...",
  "user": {
    "id": "550e8400-e29b-41d4-a716-446655440000",
    "email": "chef@gmail.com",
    "verified": true,
    "authProvider": "google"
  },
  "isNewUser": false
}
\end{lstlisting}

\subsubsection{Implementation Notes}
\begin{itemize}[noitemsep]
  \item Validera OAuth code mot Google Token API
  \item Hämta user info (email, name) från Google
  \item Kontrollera om user finns i databas (sök på email)
  \item Om finns: Logga in, generera JWT tokens, returnera \texttt{isNewUser: false}
  \item Om ny: Skapa user med \texttt{verified: true}, \texttt{auth\_provider: 'google'}, \texttt{password\_hash: null}, returnera \texttt{isNewUser: true}
  \item Generera access token (15 min expiry) och refresh token (30 dagar)
  \item Spara refresh token i \texttt{refresh\_tokens} tabell
  \item Frontend sparar tokens i localStorage och navigerar till \texttt{/inledning}
\end{itemize}

\subsection{Endpoint 3: GET /api/auth/check-email}

\subsubsection{Beskrivning}
Real-time validering om email redan är registrerad (när användare skriver i formuläret).

\subsubsection{Request}
\begin{verbatim}
GET /api/auth/check-email?email=chef@revision.se
\end{verbatim}

\subsubsection{Response - Success (200 OK)}
\textbf{Tillgänglig:}
\begin{lstlisting}
{
  "available": true
}
\end{lstlisting}

\textbf{Upptagen:}
\begin{lstlisting}
{
  "available": false,
  "message": "Email redan registrerad"
}
\end{lstlisting}

\subsubsection{Implementation Notes}
\begin{itemize}[noitemsep]
  \item Rate limit: 10 requests/minut per IP
  \item Query \texttt{users} tabell: \texttt{SELECT id FROM users WHERE email = ?}
  \item Om finns: \texttt{available: false}
  \item Används för UX (visa rött meddelande under input-fältet)
\end{itemize}

\subsection{Cloudflare Turnstile Integration}

\subsubsection{Frontend Implementation}
\begin{enumerate}[noitemsep]
  \item Ladda Cloudflare Turnstile script:
\begin{verbatim}
<script src="https://challenges.cloudflare.com/turnstile/v0/api.js" 
        async defer></script>
\end{verbatim}
  \item Rendera widget i React:
\begin{lstlisting}
<div className="cf-turnstile" 
     data-sitekey="din-site-key-har"
     data-callback="onTurnstileSuccess"></div>
\end{lstlisting}
  \item Hantera callback:
\begin{lstlisting}
const onTurnstileSuccess = (token) => {
  setCaptchaToken(token);
};
\end{lstlisting}
  \item Skicka token i \texttt{POST /api/auth/register} request
\end{enumerate}

\subsubsection{Backend Validation}
\begin{lstlisting}
import requests

def verify_turnstile(token: str, remote_ip: str) -> bool:
    response = requests.post(
        'https://challenges.cloudflare.com/turnstile/v0/siteverify',
        json={
            'secret': 'din-secret-key',
            'response': token,
            'remoteip': remote_ip
        }
    )
    result = response.json()
    return result.get('success', False)
\end{lstlisting}

\subsection{Sammanfattning}
\begin{itemize}[noitemsep]
  \item \textbf{3 API endpoints:} register, google, check-email
  \item \textbf{Bot-skydd:} Cloudflare Turnstile (gratis, proffsig UX)
  \item \textbf{Verifiering:} 6-siffrig kod via email (inte lankbaserad)
  \item \textbf{OAuth:} Google login/registrering (hoppar over email-verifiering)
  \item \textbf{External integrations:} Cloudflare API, Google OAuth API, SendGrid
\end{itemize}

\newpage

% ============================================================================
% SIDA 3: EMAIL-VERIFIERING (Autentiseringssida utan sidebar)
% ============================================================================

\section{Email-verifiering (Sida 3)}

\textit{TODO: Dokumentera verifierings-endpoints här...}

\newpage

% ============================================================================
% SIDA 4: LOGIN (Autentiseringssida utan sidebar)
% ============================================================================

\section{Login (Sida 4)}

\subsection{Översikt}
Inloggningssidan låter användare logga in med e-post och lösenord, eller via Google OAuth. Sidan innehåller även länk till ``Glömt lösenord?'' som leder till lösenordsåterställning. Vid lyckad inloggning får användaren JWT-tokens som sparas i localStorage.

\subsection{Funktionalitet}
\begin{itemize}[noitemsep]
  \item Användaren anger e-post och lösenord
  \item Vid ``Logga in'' görs POST till \texttt{/api/auth/login}
  \item Alternativt klickar användaren på ``Logga in med Google'' (OAuth-flöde)
  \item Vid lyckad inloggning:
  \begin{itemize}
    \item Backend returnerar \texttt{accessToken} (15 min) och \texttt{refreshToken} (30 dagar)
    \item Frontend sparar tokens i \texttt{localStorage}
    \item Användaren navigeras till \texttt{/inledning} (första slide med sidebar)
  \end{itemize}
  \item Vid fel visas felmeddelande (``Fel e-post eller lösenord'', ``Kontot är inte verifierat'')
  \item Länk ``Glömt lösenord?'' navigerar till \texttt{/forgot-password}
\end{itemize}

\subsection{API Endpoints}

\subsubsection{POST /api/auth/login}
\textbf{Beskrivning:} Loggar in användare med e-post och lösenord.

\textbf{Request body (JSON):}
\begin{lstlisting}[language=json]
{
  "email": "user@example.com",
  "password": "MySecurePass123"
}
\end{lstlisting}

\textbf{Validering:}
\begin{itemize}[noitemsep]
  \item E-post finns i databasen (annars 401 \texttt{invalid\_credentials})
  \item Lösenord matchar bcrypt-hash (annars 401 \texttt{invalid\_credentials})
  \item Användaren är verifierad (\texttt{verified = true}, annars 403 \texttt{unverified\_email})
\end{itemize}

\textbf{Success Response (200 OK):}
\begin{lstlisting}[language=json]
{
  "accessToken": "eyJhbGciOiJIUzI1NiIsInR5cCI6IkpXVCJ9...",
  "refreshToken": "eyJhbGciOiJIUzI1NiIsInR5cCI6IkpXVCJ9...",
  "user": {
    "id": "uuid-123",
    "email": "user@example.com",
    "role": "user",
    "firmId": "firm-uuid-456"
  }
}
\end{lstlisting}

\textbf{Error Responses:}
\begin{itemize}[noitemsep]
  \item \textbf{400 Bad Request:} \texttt{\{"error": "missing\_fields", "message": "Email and password required"\}}
  \item \textbf{401 Unauthorized:} \texttt{\{"error": "invalid\_credentials", "message": "Incorrect email or password"\}}
  \item \textbf{403 Forbidden:} \texttt{\{"error": "unverified\_email", "message": "Please verify your email first"\}}
  \item \textbf{429 Too Many Requests:} Rate limit (5 req/min per IP)
\end{itemize}

\textbf{Implementation notes:}
\begin{itemize}[noitemsep]
  \item Använd \texttt{bcrypt.compare()} för lösenordskontroll
  \item Generera JWT med \texttt{jwt.sign()} (HS256 algoritm)
  \item Spara \texttt{refreshToken} i \texttt{refresh\_tokens} tabell med \texttt{expires\_at}
  \item Rate limiting: 5 försök per minut per IP-adress (förhindra brute force)
  \item Logga misslyckade inloggningsförsök för säkerhet
\end{itemize}

\textbf{Database queries:}
\begin{verbatim}
SELECT id, email, password_hash, verified, role, firm_id 
FROM users 
WHERE email = ?;

INSERT INTO refresh_tokens (user_id, token, expires_at) 
VALUES (?, ?, NOW() + INTERVAL '30 days');
\end{verbatim}

\subsubsection{POST /api/auth/google}
\textbf{Beskrivning:} Loggar in eller registrerar användare via Google OAuth.

\textbf{OBS:} Denna endpoint används även i Section 3 (Registrering). Samma implementation används för både login och registrering eftersom Google OAuth hanterar båda fallen automatiskt.

\textbf{Request body (JSON):}
\begin{lstlisting}[language=json]
{
  "code": "4/0AY0e-g7X...",
  "redirectUri": "http://localhost:5173/login"
}
\end{lstlisting}

\textbf{Flöde:}
\begin{enumerate}[noitemsep]
  \item Frontend öppnar Google OAuth popup med \texttt{client\_id}
  \item Användaren loggar in och godkänner åtkomst
  \item Google returnerar \texttt{code} till frontend
  \item Frontend skickar \texttt{code} till backend
  \item Backend byter \texttt{code} mot \texttt{access\_token} via Google API
  \item Backend hämtar användarinfo från Google (email, name, picture)
  \item Om användaren finns: Logga in (returnera JWT tokens)
  \item Om användaren INTE finns: Skapa ny användare med \texttt{verified=true} och \texttt{auth\_provider='google'}
\end{enumerate}

\textbf{Success Response (200 OK för login, 201 Created för ny användare):}
\begin{lstlisting}[language=json]
{
  "accessToken": "eyJhbGciOiJIUzI1NiIsInR5cCI6IkpXVCJ9...",
  "refreshToken": "eyJhbGciOiJIUzI1NiIsInR5cCI6IkpXVCJ9...",
  "user": {
    "id": "uuid-123",
    "email": "user@gmail.com",
    "name": "John Doe",
    "verified": true,
    "authProvider": "google"
  },
  "isNewUser": false
}
\end{lstlisting}

\textbf{Error Responses:}
\begin{itemize}[noitemsep]
  \item \textbf{400 Bad Request:} \texttt{\{"error": "invalid\_code", "message": "Invalid or expired authorization code"\}}
  \item \textbf{503 Service Unavailable:} \texttt{\{"error": "google\_api\_error", "message": "Google authentication temporarily unavailable"\}}
\end{itemize}

\textbf{External API call:}
\begin{verbatim}
POST https://oauth2.googleapis.com/token
{
  "code": "...",
  "client_id": "...",
  "client_secret": "...",
  "redirect_uri": "...",
  "grant_type": "authorization_code"
}

GET https://www.googleapis.com/oauth2/v2/userinfo
Headers: { Authorization: "Bearer <access_token>" }
\end{verbatim}

\subsection{Frontend Implementation}

\textbf{LoginSlide.jsx - Email/Password login:}
\begin{lstlisting}[language=javascript]
const handleLogin = async (e) => {
  e.preventDefault();
  try {
    const response = await fetch('http://localhost:8000/api/auth/login', {
      method: 'POST',
      headers: { 'Content-Type': 'application/json' },
      body: JSON.stringify({ email, password })
    });
    
    if (!response.ok) {
      const error = await response.json();
      if (error.error === 'unverified_email') {
        setError('Verifiera din e-post forst');
      } else {
        setError('Fel e-post eller losenord');
      }
      return;
    }
    
    const data = await response.json();
    localStorage.setItem('accessToken', data.accessToken);
    localStorage.setItem('refreshToken', data.refreshToken);
    onNext(); // Navigate to /inledning
  } catch (err) {
    setError('Något gick fel. Försök igen.');
  }
};
\end{lstlisting}

\textbf{LoginSlide.jsx - Google OAuth login:}
\begin{lstlisting}[language=javascript]
const handleGoogleLogin = () => {
  const clientId = 'YOUR_GOOGLE_CLIENT_ID';
  const redirectUri = 'http://localhost:5173/login';
  const scope = 'email profile';
  
  const authUrl = 
    `https://accounts.google.com/o/oauth2/v2/auth?` +
    `client_id=${clientId}&` +
    `redirect_uri=${redirectUri}&` +
    `response_type=code&` +
    `scope=${scope}`;
  
  window.location.href = authUrl;
};

// Efter redirect tillbaka från Google:
useEffect(() => {
  const urlParams = new URLSearchParams(window.location.search);
  const code = urlParams.get('code');
  
  if (code) {
    fetch('http://localhost:8000/api/auth/google', {
      method: 'POST',
      headers: { 'Content-Type': 'application/json' },
      body: JSON.stringify({ code, redirectUri: window.location.origin + '/login' })
    })
    .then(res => res.json())
    .then(data => {
      localStorage.setItem('accessToken', data.accessToken);
      localStorage.setItem('refreshToken', data.refreshToken);
      navigate('/inledning');
    });
  }
}, []);
\end{lstlisting}

\subsection{Database Tables}

\textbf{users:}
\begin{verbatim}
id              UUID PRIMARY KEY
email           VARCHAR(255) UNIQUE NOT NULL
password_hash   VARCHAR(255) -- NULL för Google OAuth-användare
verified        BOOLEAN DEFAULT FALSE
role            VARCHAR(50) DEFAULT 'user'
firm_id         UUID REFERENCES firms(id)
auth_provider   VARCHAR(50) DEFAULT 'email' -- 'email' eller 'google'
created_at      TIMESTAMP DEFAULT NOW()
\end{verbatim}

\textbf{refresh\_tokens:}
\begin{verbatim}
id          UUID PRIMARY KEY
user_id     UUID REFERENCES users(id)
token       TEXT NOT NULL
expires_at  TIMESTAMP NOT NULL
created_at  TIMESTAMP DEFAULT NOW()
\end{verbatim}

\subsection{Sammanfattning - Login endpoints}
\begin{itemize}[noitemsep]
  \item \textbf{POST /api/auth/login} -- Logga in med email/password, returnera JWT tokens
  \item \textbf{POST /api/auth/google} -- Google OAuth login/registrering (samma som Section 3)
  \item \textbf{Glömt lösenord-länk} -- Navigerar till \texttt{/forgot-password} (se Section 4.5)
  \item Rate limiting: 5 försök per minut per IP-adress
  \item Tokens sparas i localStorage: \texttt{accessToken} (15 min), \texttt{refreshToken} (30 dagar)
\end{itemize}

\newpage

% ############################################################################
% ############################################################################
% DEL 2: ONBOARDING SLIDES (MED SIDEBAR)
% ############################################################################
% ############################################################################

% ============================================================================
% SLIDE 1: INTRO (Första slide med sidebar)
% ============================================================================

\section{Intro (Slide 5)}

\textit{TODO: Dokumentera intro-endpoints här (första slide med sidebar)...}

\newpage

% ============================================================================
% SLIDE 2: UPPDRAG
% ============================================================================

\section{Uppdrag (Slide 6)}

\textit{TODO: Dokumentera uppdrag-endpoints här...}

\newpage

% ============================================================================
% SLIDE 3: VERKSAMHET
% ============================================================================

\section{Verksamhet (Slide 7)}

\textit{TODO: Dokumentera verksamhet-endpoints här...}

\newpage

% ============================================================================
% SLIDE 4: RISKFRÅGOR
% ============================================================================

\section{Riskfrågor (Slide 8)}

\textit{TODO: Dokumentera riskfrågor-endpoints här...}

\newpage

% ============================================================================
% SLIDE 5: IDENTITETSKONTROLL
% ============================================================================

\section{Identitetskontroll (Slide 9)}

\textit{TODO: Dokumentera BankID-endpoints här...}

\newpage

\section{Kontrolltabell (Slide 10)}

\textit{TODO: Dokumentera kontrolltabell-endpoints här...}

\newpage

% ============================================================================
% SLIDE 7: PEP-FÖRDJUPNING
% ============================================================================

\section{PEP-fördjupning (Slide 11)}

\textit{TODO: Dokumentera PEP-endpoints här...}

\newpage

% ============================================================================
% SLIDE 8: VERKSAMHETSKOD
% ============================================================================

\section{Verksamhetskod (Slide 12)}

\textit{TODO: Dokumentera SNI-kod-endpoints här...}

\newpage

% ============================================================================
% SLIDE 9: ÄGARSTRUKTUR
% ============================================================================

\section{Ägarstruktur (Slide 13)}

\textit{TODO: Dokumentera Bolagsverket-endpoints här...}

\newpage

% ============================================================================
% SLIDE 10: STYRELSE
% ============================================================================

\section{Styrelse (Slide 14)}

\textit{TODO: Dokumentera styrelse-endpoints här...}

\newpage

% ============================================================================
% SLIDE 11: FÖRETAGSDOKUMENTATION
% ============================================================================

\section{Företagsdokumentation (Slide 13.5)}

\textit{TODO: Dokumentera PDF-upload och OCR-endpoints här...}

\newpage

% ============================================================================
% SLIDE 12: BOKFÖRINGSUNDERLAG
% ============================================================================

\section{Bokföringsunderlag (Slide 14.5)}

\textit{TODO: Dokumentera file-upload och fraud-detection-endpoints här...}

\newpage

% ============================================================================
% SLIDE 13: BOKFÖRING
% ============================================================================

\section{Bokföring (Slide 15)}

\textit{TODO: Dokumentera Skatteverket/SIE-endpoints här...}

\newpage

% ============================================================================
% SLIDE 14: LIKVIDITET
% ============================================================================

\section{Likviditet (Slide 16)}

\textit{TODO: Dokumentera likviditets-analys-endpoints här...}

\newpage

% ============================================================================
% SLIDE 15: OMSÄTTNING
% ============================================================================

\section{Omsättning (Slide 17)}

\textit{TODO: Dokumentera omsättnings-data-endpoints här...}

\newpage

% ============================================================================
% SLIDE 16: RESULTAT
% ============================================================================

\section{Resultat (Slide 18)}

\textit{TODO: Dokumentera resultat-analys-endpoints här...}

\newpage

% ============================================================================
% SLIDE 17: BRANSCH
% ============================================================================

\section{Bransch (Slide 19)}

\textit{TODO: Dokumentera bransch-jämförelse-endpoints här...}

\newpage

% ============================================================================
% SLIDE 18: BOKFÖRINGSANALYS
% ============================================================================

\section{Bokföringsanalys (Slide 20)}

\textit{TODO: Dokumentera AI-analys-endpoints här...}

\newpage

% ============================================================================
% SLIDE 19: RISKBESKED
% ============================================================================

\section{Riskbesked (Slide 21)}

\textit{TODO: Dokumentera risk-scoring-endpoints här...}

\newpage

% ============================================================================
% SLIDE 20: SKYLDIGHETER
% ============================================================================

\section{Skyldigheter (Slide 22)}

\textit{TODO: Dokumentera avtal-endpoints här...}

\newpage

% ============================================================================
% SLIDE 21: FORTNOX-KOPPLING
% ============================================================================

\section{Fortnox-koppling (Slide 23)}

\textit{TODO: Dokumentera Fortnox OAuth-endpoints här...}

\newpage

% ============================================================================
% SLIDE 22: BANKRÄTTIGHETER
% ============================================================================

\section{Bankrättigheter (Slide 24)}

\textit{TODO: Dokumentera bank-OAuth-endpoints här...}

\newpage

% ############################################################################
% ############################################################################
% DEL 3: ADMIN DASHBOARD (Separat spec finns i Admin_Dashboard_Spec.tex)
% ############################################################################
% ############################################################################

\section{AdminDashboard}

\textit{TODO: Dokumentera admin-endpoints här (referera till Admin\_Dashboard\_Spec.tex)...}

\newpage

\section{Appendix}

\subsection{Gemensamma Error Codes}
\begin{table}[h]
\centering
\begin{tabular}{@{}lll@{}}
\toprule
\textbf{Code} & \textbf{Status} & \textbf{Beskrivning} \\ \midrule
400 & Bad Request & Ogiltig input data \\
401 & Unauthorized & Saknar eller ogiltig JWT token \\
403 & Forbidden & Saknar behörighet för resursen \\
404 & Not Found & Resursen finns inte \\
409 & Conflict & Resurs existerar redan (t.ex. email) \\
422 & Unprocessable Entity & Valideringsfel \\
429 & Too Many Requests & Rate limit överskriden \\
500 & Internal Server Error & Oväntat serverfel \\
503 & Service Unavailable & Extern tjänst ej tillgänglig \\ \bottomrule
\end{tabular}
\end{table}

\subsection{Rate Limiting}
\begin{itemize}[noitemsep]
  \item \textbf{Auth endpoints:} 5 requests/minut per IP
  \item \textbf{File uploads:} 10 requests/minut per user
  \item \textbf{Övriga endpoints:} 100 requests/minut per user
\end{itemize}

\subsection{File Upload Limits}
\begin{itemize}[noitemsep]
  \item \textbf{Max filstorlek:} 10 MB per fil
  \item \textbf{Max totalt per request:} 50 MB
  \item \textbf{Tillåtna format:} PDF, PNG, JPG, JPEG, SIE, XML
\end{itemize}

\end{document}
